% 文档与参赛身份配置。
% 这里定义全项目统一占位值，必须最先载入。
% 未确认字段保留 \PendingValue；它表示尚待审阅，不能视为承诺或成果。
\newcommand{\PendingValue}{待人工审阅}

% 文档管理字段。版本号表示排版草案版本，不代表技术方案已完成。
\newcommand{\DocumentTitle}{项目计划书暨命题对策书}
\newcommand{\DocumentSubtitle}{排版与配置草案}
\newcommand{\DocumentVersion}{v0.1}
\newcommand{\DocumentDate}{2026年9月8日}
\newcommand{\DocumentStatus}{排版草案，技术内容待人工审阅}
\newcommand{\DocumentAuthor}{\PendingValue}
\newcommand{\DocumentReviewer}{\PendingValue}

% 已知比赛和企业命题信息；学校与具体子方向仍由团队核定。
\newcommand{\CompetitionName}{中国国际大学生创新大赛（2026）}
\newcommand{\CompetitionTrack}{产业赛道}
\newcommand{\CompetitionGroup}{企业命题组}
\newcommand{\CompetitionDirection}{\PendingValue}
\newcommand{\SchoolName}{\PendingValue}
\newcommand{\CollegeName}{\PendingValue}
% 两段组合为完整题名，封面在语义边界换行，正文使用完整题名宏。
\newcommand{\ChallengeTitleLineOne}{基于国产开源软件与开源RISC-V处理器架构的}
\newcommand{\ChallengeTitleLineTwo}{嵌入式智能计算系统研究与设计}
\newcommand{\ChallengeTitle}{\ChallengeTitleLineOne\ChallengeTitleLineTwo}
\newcommand{\ChallengeCompany}{中智讯（武汉）科技有限公司}
\newcommand{\ChallengeIdentifier}{239}
\newcommand{\ProjectName}{\PendingValue}
\newcommand{\ProjectSubtitle}{\PendingValue}
\newcommand{\ApplicationScenario}{\PendingValue}
\newcommand{\SubmissionStage}{校内选拔赛材料准备}
\newcommand{\SubmissionDeadline}{2026年9月13日12时}
